\section{总结与延伸}

\subsection{总结表}

\begin{table}[H]
\centering
\begin{tabular}{p{0.2\textwidth} p{0.26\textwidth} p{0.26\textwidth} p{0.16\textwidth}}
\toprule
领域 & 核心策略 & 实施手段 & 进一步问题 \\
\midrule
Benchmarking & 场景化评估＋rubric & Decision support case + POAP + slide rubric & 如何量化 alignment drift？ \\
Monitoring & Telemetry + replay + alerts & Metrics/alert logs + incident loop & 怎么在多 agent 运行时保持解释性？ \\
Governance & Pre-deploy gates + Trust oversight & red teams + capability reports + trustee veto & 怎样把 governance artifact 嵌入 CI/CD？ \\
Research & Capability contour + multi-agent orchestration & Long-term Benefit Trust + measurement plan & 如何抓住 tail-end risks before deployment？ \\
\bottomrule
\end{tabular}
\caption{Lecture 02 的能力、安全、治理、研究四层面压缩表}
\end{table}

\subsection{进一步阅读}

\begin{itemize}
    \item Anthropic, ``Responsible Scaling Policy'' 2023
    \item Anthropic, ``Claude Model Cards and Capability Reports'' 2024
    \item Anthropic, ``Long-term Benefit Trust'' governance brief
    \item Ben Mann, talk materials from Berkeley RDI (slides included in repo)
    \item Shevlane et al., ``Model Evaluation for Extreme Risks,'' 2023
\end{itemize}

\subsection{本章小结}

Lecture 02 把能力评估、监控、部署 guardrail 和制度治理串成一条闭环：不断评估 + 监控 + pre-deploy gates + trustee oversight；研究议程则将 benchmark measurement、policy orchestration 与 governance artifact 进一步向前拉，为未来的 Agent 安全奠定行动框架。

\subsection{Slide Evidence}

The lecture slides underline the emphasis on evidence-backed governance. Slide 4 distills the capability matrix, while slide 5 captures the Trust handoff timeline.

\begin{figure}[H]
\centering
\includegraphics[page=4,width=0.92\textwidth]{slides.pdf}
\caption{Slide 4: Capability matrix and risk signals.\protect\footnotemark}
\end{figure}

\begin{figure}[H]
\centering
\includegraphics[page=5,width=0.92\textwidth]{slides.pdf}
\caption{Slide 5: Long-term Benefit Trust timeline.\protect\footnotemark}
\end{figure}
\footnotetext{Screenshots extracted from the provided slide deck.}

\end{document}
